\section{Limitations and Future work}
\label{sec:limitation}
One of \Ansor's limitations is that \Ansor cannot optimize graphs with dynamic shapes~\cite{shen2020nimble}.
\Ansor requires the shapes in the computational graph to be static and known in advance
to do analysis, construct the search space, and perform measurements. How to generate programs for symbolic or dynamic shape is an interesting future direction.
Another limitation is that \Ansor only supports dense operators. 
To support sparse operators (\textit{e.g.}, SpMM) that are commonly used in sparse neural networks~\cite{gale2019state} and graph neural networks~\cite{hu2020featgraph}, we expect that a large portion of \Ansor can still be reused, but we need to redesign the search space.
Lastly, \Ansor only performs program optimizations at a high level but relies on other code generators (\textit{e.g.}, LLVM and NVCC) to do platform-dependent optimizations ({\textit{e.g.}, instruction selection}). \Ansor comes short of utilizing the special instructions, such as Intel VNNI, NVIDIA Tensor Core, and ARM Dot for mixed-precision and low-precision operators, which are not handled well by the off-the-shelf code generators currently.

\section{Conclusion}
\label{sec:conclusion}

We propose \Ansor, an automated search framework that generates high-performance tensor programs for deep neural networks. 
By efficiently exploring a large search space and prioritizing performance bottlenecks,
\Ansor finds high-performance programs that are outside the search space of existing approaches.
\Ansor outperforms existing manual libraries and search-based frameworks on a diverse set of neural networks and hardware platforms by up to $3.8 \times$.
By automatically searching for better programs, we hope that \Ansor will help bridge the gap between the increasing demand in computing power and limited hardware performance.
\Ansor is integrated into the Apache TVM open-source project \footnote{\url{https://tvm.apache.org/}}.
